Where people commute defines their economic community more precisely than
where they work or where they live considered in isolation.
Two residential census tracts whose workers flow to the same employment
centers share a functional economic area: they patronize the same
commercial corridors, depend on the same transportation infrastructure,
and have a shared stake in the policy decisions that govern their common
destination.
This paper presents M.4, the fourth paper in the M-track community character
series, which operationalizes commuting community membership through the
Census Bureau's LODES Origin-Destination (OD) file.

For each census tract we compute a commuting destination distribution
$D(h) = \{w : \mathrm{count}(h \to w)\}$: a sparse vector recording how many
workers residing in tract $h$ are employed in each destination tract $w$.
The community similarity between two home tracts is the weighted Jaccard
similarity of their destination distributions: the overlap of commute flows
divided by the union of commute flows.
Tracts whose workers commute to the same employment centers receive high
similarity and heavy edge weights; tracts serving different economic spheres
receive low similarity and light edges that invite district cuts at functional
economic area boundaries.

Experimental design targets Wisconsin ($k=8$) and North Carolina ($k=14$)
with primary metrics of suburban clustering rates and polycentric separation.
The Wisconsin prediction---Milwaukee suburb tracts sharing a primary
destination co-assign at ${\geq}70\%$ rate---and the North Carolina
prediction---Raleigh metro and Charlotte metro tracts achieve cross-metro
co-assignment below 5\% under commuting-shed weights---are both deferred
to Phase~2 pending LODES OD data download.\textsuperscript{$\dagger$}
The LODES OD data is annual, free, publicly available, and contains no
racial or partisan content.
